\section{Supplementary Reviewer Guide}

This supplement is organized as an explanatory companion for the simulation-checked diagnostic story in
the main paper. It does not add a second paper hidden behind the page limit. Instead, it gives the
definitions, equations, implementation notes, and new visual explanations that make the main
paper easier to audit. The plates show predicate anatomy, generated-package consumption, link
ownership, failure labels, and provenance flow while keeping the main paper as the complete
statement of the contribution and evidence scope. The supplement is a slower reading path for
reviewers who want to inspect how a candidate primitive package becomes a Newton task artifact.

\paragraph{How to read this supplement.}
Sections~2--5 are the technical core. Section~2 defines the package model and the body/link
attachment map. Section~3 derives the diagnostic predicates used by the small probe scenes.
Section~4 explains why compound body state can change when locally plausible primitive shapes are
assembled into a single simulated body. Section~5 explains the robot-specific link ownership rule.
Sections~6--8 then collect new visual atlas pages, reproducibility notes, and claim-boundary
checks. The sequence is meant to be read independently of the main-paper figure order.

\paragraph{What is new here.}
The supplement figures are new explanatory figures: predicate diagrams, failure storyboards,
Franka link ownership sketches, source-shape suppression accounting, candidate-lane anatomy, and
provenance flow. They are not copies of main-paper figures. The supplement tables are notation,
predicate, parameter, provenance, and boundary tables. They are not copies of main-paper tables and
do not restate the main quantitative result table.

\paragraph{What remains bounded.}
The wording follows the same claim registry as the main paper. A package can be called
simulation-checked only when a dated record links it to a named Newton task, settings, asset,
environment, and report. The diagnostic checker is a runtime acceptance layer for scoped tasks, not
a formal verification system. The Franka examples teach link ownership and generated-package
consumption for a recorded smoke asset; they should not be read as manipulation, whole-arm
operation, deployment, real-world transfer, or safety certification evidence.

\begin{table}[t]
\centering
\small
\caption{Reading map for the supplement. The entries identify what each section adds beyond the
main text without expanding the claim boundary.}
\label{tab:supp-reading-map}
\begin{tabular}{p{0.18\linewidth}p{0.35\linewidth}p{0.34\linewidth}}
\toprule
Section & Reviewer question & Supplement role \\
\midrule
S2 & What exactly is a primitive package? & Defines symbols, ownership, body attachment, and outcomes. \\
S3 & How do probe labels arise? & Derives predicate clauses and explains threshold accounting. \\
S4 & Why can a good-looking primitive fail? & Explains aggregate COM and inertia through the parallel-axis theorem. \\
S5 & What is robot-specific? & Defines link ownership, cross-link merge rejection, and source-shape suppression. \\
S6 & What should I look for visually? & Provides explanatory plates and failure storyboards. \\
S7 & How are artifacts traced? & Links configs, records, figure manifests, and paper assets. \\
S8 & What should I not infer? & Restates unsupported claims and the evidence needed to strengthen them. \\
\bottomrule
\end{tabular}
\end{table}

The rest of the supplement uses compact reading notes after figures. Each note names the intended
inspection target and the unsupported inference to avoid. This is deliberate: the paper studies
safety-affecting collision packages, so visual clarity is useful only when paired with explicit
evidence limits.
